\section{Methodology}
\label{sec:method}

We ask four research questions.
\textbf{RQ1}: when real open-weight models are attacked, how often do they comply, how often is harm
achieved, and what fraction of compromised runs does the pack detect?
\textbf{RQ2}: on public datasets of attacks and of real benign traffic, what recall and false-positive
rate do the content signatures achieve?
\textbf{RQ3}: how do they compare with, and combine with, open ML prompt-injection classifiers?
\textbf{RQ4}: how does detection degrade under an adaptive attacker who mutates the payload, and how
many mutants evade detection while the attack still works?

\subsection{Rule execution}
All rules run unchanged as KQL on the Kusto emulator (Kustainer, image pinned by digest). The only
modification is to widen a \texttt{trigger\_window} parameter that deduplicates alerts across
overlapping scheduled runs, because replayed data is not ``recent''. Events are loaded into a table
with the production schema (\texttt{LLMAppLogs\_CL}).

\subsection{RQ1: live attacks}
\paragraph{Models.} Five open-weight instruction models from five vendors, all with tool calling and
small enough for an 8~GB consumer GPU (NVIDIA RTX~4060), served by Ollama~0.40.6 at default
quantisation: \texttt{gemma4:e4b}, \texttt{granite4:tiny-h}, \texttt{llama3.1:8b}, \texttt{ministral-3:8b}, \texttt{qwen3:8b}.

\paragraph{Scenarios.} Ten attack scenarios in which only attacker inputs are fixed and the model's
behaviour is not scripted: a direct injection asking for the system prompt; indirect injections in
two poisoned vendor documents (email exfiltration, image beacon); MCP tool poisoning, rug pull and
runtime-response poisoning against a real MCP server; a two-session memory poisoning; agent
configuration tampering and a reverse shell delivered through pasted text; and a base64-obfuscated
injection. The system prompt instructs the assistant not to send email unless asked, and no
high-risk tool requires approval, which is a common and insecure default.

\paragraph{Oracles.} Each scenario has two independent oracles defined on side effects, never on the
rules: \emph{compliance} (the model took the attacker-directed step, for example it read the
restricted file the poisoned tool description asked for) and \emph{harm} (the attacker's goal was
achieved, for example restricted data reached the third-party tool). Oracles are unit-tested on
positive and negative traces. Each (model, scenario) pair runs 10 times at temperature 0.7 with fixed
seeds, on a fresh agent instance. Fifty-eight benign prompts (8 application questions and 50 real
prompts drawn from OASST2 and Dolly) measure live false positives.

\paragraph{Metrics.} Compliance and harm rates; detection given compliance (one of the scenario's
intended rules fired on the run); detection given harm; any-rule visibility of resisted attempts; and
the false-positive rate on benign runs. Proportions carry Wilson 95\% intervals.

\subsection{RQ2: public datasets}
We assembled 112{,}894 items from 16 public datasets (Table~\ref{tab:rq2-datasets}): labelled direct
injections and jailbreaks (deepset, Lakera Gandalf, neuralchemy, SPML, jackhhao, TrustAIRLab
in-the-wild), indirect injections embedded in documents (Microsoft's LLMail-Inject challenge
submissions~\cite{llmail}, BIPIA~\cite{bipia}, NVIDIA's agentic indirect injection set), poisoned MCP
tool descriptions from MCPTox~\cite{mcptox}, and benign material: real user prompts (WildChat-1M,
OASST2), crowd-written instructions and contexts (Dolly-15k), and real MCP tool descriptions
(MCP-Flow). Prompts become \texttt{prompt} events, documents become \texttt{retrieval} events and
tool results, and tool descriptions become \texttt{tool\_registry} events. WildChat is real traffic
that can itself contain attacks, so it is reported as an alert rate rather than a false-positive rate.
Intervals are percentile bootstrap (1{,}000 resamples).

\subsection{RQ3: ML classifiers}
Two open classifiers run with ONNX Runtime on CPU: ProtectAI's DeBERTa-v3 prompt-injection model
(v2)~\cite{protectai} and Meta's Llama Prompt Guard~2 (22M)~\cite{promptguard}, both at threshold 0.5.
ProtectAI's model card names jackhhao among its training data and lists eight unnamed MIT-licensed
sets, so we mark jackhhao as \emph{contaminated} and the MIT-licensed sets published before the model
(Lakera Gandalf, SPML, TrustAIRLab) as \emph{possibly contaminated}; headline numbers use the
remaining, clean datasets only. We also evaluate the union and intersection of rules and classifier,
and a precision-recall sweep.

\subsection{RQ4: adaptive attacker}
Twelve mutation operators: deterministic synonym substitution, homoglyphs, leetspeak, zero-width
insertion, word splitting and base64 wrapping; and LLM-based paraphrase and translation into German,
Hungarian, Spanish, Chinese and Russian (Llama~3.1~8B, temperature~0, fixed seed). \emph{Offline},
we mutate 300 real attack prompts that the rules detect at baseline and re-run rules and classifiers.
\emph{Live}, we mutate the attacker-controlled text of the direct (L01) and indirect (L02) scenarios
and re-run them on two models, counting \emph{evasion-preserving} mutants: the model complied and no
intended rule fired.
